%! Author = trdng
%! Date = 4/4/2026

\documentclass[conference]{IEEEtran}

\usepackage{amsmath, amssymb, braket}

\title{Quantum Fourier Transform and Its Circuit Construction}

\author{
\IEEEauthorblockN{}
\IEEEauthorblockA{}
}

\begin{document}

\maketitle

\begin{abstract}
The Quantum Fourier Transform (QFT) is a fundamental primitive in quantum computing, forming the basis of several important algorithms such as phase estimation, Shor's factoring algorithm, and hidden subgroup problems. This paper presents the formal definition of QFT, its binary decomposition, and an efficient quantum circuit construction.
\end{abstract}

\section{Introduction}\label{sec:introduction}

The Quantum Fourier Transform is the quantum analogue of the discrete Fourier transform (DFT). It operates on quantum states and enables exponential speedups in several quantum algorithms.

\section{Definition of QFT}\label{sec:definition-of-qft}

Let $N = 2^n$. The QFT is defined as a unitary transformation acting on computational basis states:

\begin{equation}
\mathrm{QFT}\ket{j} = \frac{1}{\sqrt{N}} \sum_{k=0}^{N-1} e^{\frac{2\pi i j k}{N}} \ket{k}\label{eq:equation}
\end{equation}

Where $j \in \{0, \dots, N-1\}$.

The QFT operator is unitary, satisfying:
\begin{equation}
\langle Uj | Uj' \rangle = \delta_{jj'}\label{eq:equation2}
\end{equation}

\section{Binary Decomposition}

Let the binary representation of $j$ be:
\begin{equation}
j = \sum_{x=1}^{n} j_x 2^{n-x}, \quad j_x \in \{0,1\}\label{eq:equation3}
\end{equation}

Then the QFT can be expressed in tensor product form:
\begin{equation}
\mathrm{QFT}\ket{j}
= \bigotimes_{x=1}^{n}
\frac{1}{\sqrt{2}}
\left(
\ket{0} + e^{2\pi i \, 0.j_x j_{x+1} \dots j_n} \ket{1}
\right)\label{eq:equation4}
\end{equation}

where:
\begin{equation}
0.j_x j_{x+1}\dots j_n =
\frac{j_x}{2} + \frac{j_{x+1}}{2^2} + \cdots + \frac{j_n}{2^{n-x+1}}\label{eq:equation5}
\end{equation}

\section{Quantum Circuit Construction}\label{sec:quantum-circuit-construction}

The QFT circuit is constructed using two fundamental gates:

\subsection{Hadamard Gate}\label{subsec:hadamard-gate}

\begin{equation}
H\ket{j} =
\frac{1}{\sqrt{2}}
\left(
\ket{0} + e^{2\pi i \, 0.j} \ket{1}
\right)\label{eq:equation6}
\end{equation}

\subsection{Controlled Phase Rotation}\label{subsec:controlled-phase-rotation}

Define the controlled rotation gate $CR_k$ with phase:
\begin{equation}
R_k =
\begin{pmatrix}
1 & 0 \\
0 & e^{\frac{2\pi i}{2^k}}
\end{pmatrix}\label{eq:equation7}
\end{equation}

\subsection{Circuit Structure}\label{subsec:circuit-structure}

The QFT circuit applies:

\begin{itemize}
    \item A Hadamard gate on qubit $j_x$
    \item Controlled rotations $CR_k$ with control qubits $j_{x+1}, \dots, j_n$
\end{itemize}

Formally, for each qubit $j_x$:
\begin{equation}
\ket{j_x} \rightarrow
\frac{1}{\sqrt{2}}
\left(
\ket{0} + e^{2\pi i \, 0.j_x \dots j_n} \ket{1}
\right)\label{eq:equation8}
\end{equation}

Due to reversed output ordering, a sequence of SWAP gates is applied at the end of the circuit.

\section{Inverse QFT}\label{sec:inverse-qft}

The inverse QFT, denoted $\mathrm{QFT}^\dagger$, is obtained by:

\begin{itemize}
    \item Reversing the order of gates
    \item Conjugating phase rotations ($e^{i\theta} \rightarrow e^{-i\theta}$)
\end{itemize}

\section{Complexity Analysis}\label{sec:complexity-analysis}

The QFT circuit requires:
\begin{equation}
O(n^2)\label{eq:equation9}
\end{equation}

More precisely:
\begin{itemize}
    \item $\frac{n(n+1)}{2}$ Controlled rotation gates
    \item $\frac{n}{2}$ SWAP gates
\end{itemize}

\section{Conclusion}\label{sec:conclusion}

The QFT provides an efficient quantum implementation of the Fourier transform with polynomial circuit complexity. Its structured decomposition enables practical realization on quantum hardware and serves as a key building block for many quantum algorithms.

\end{document}